\documentclass[12pt]{article}

% Paquetes esenciales
\usepackage[utf8]{inputenc}
\usepackage[T1]{fontenc}
\usepackage{amsmath, amssymb}
\usepackage{graphicx}
\usepackage{booktabs}
\usepackage{caption}
\usepackage{hyperref}
\usepackage[margin=1in]{geometry}

% -------------------------------------------------------
% Configuración para numerar figuras y tablas con "S"
% -------------------------------------------------------
\renewcommand{\thefigure}{S\arabic{figure}}
\renewcommand{\thetable}{S\arabic{table}}
\renewcommand{\theequation}{S\arabic{equation}}
\renewcommand{\thesection}{S\arabic{section}}

\begin{document}
	
	% Encabezado
	\begin{center}
		{\LARGE \textbf{Supplementary Material}}\\[0.5em]
		{\large Science in the Colombian Public Sphere:  \\
			A Digital Humanities and NLP Analysis of Opinion Columns (2018–2020)}
	\end{center}
	
	\vspace{1em}
	\hrule
	\vspace{1em}
	


	\section*{Supplementary Material}
	
	Below are some excerpts that the developed pipeline classifies as belonging to the field of science, presented in their original language.

	
	From the category Science and research management
	
	\begin{quote}
		(...) separado. El Ministerio de Ciencia, Tecnología e Innovación tiene la responsabilidad de convocar la constitución de un Sistema Nacional de Innovación que asuma como propósito desarrollar la hoja de ruta planteada por la Misión de Sabios, alineando a todos los actores a nivel local, regional y nacional. El conocimiento siempre es un logro colectivo. Frente a los desafíos que afrontamos como país, es necesario que asumamos colectivamente, todos los actores, el reto de gestionar el conocimiento, aprendiendo de manera continua, para ponerlo al servicio de las comunidades. Para superar la crisis y lograr las transformaciones estructurales se requiere planeación y un trabajo de mediano y largo plazo, en donde todos sumemos sin pensar solo en réditos inmediatos. Un sistema en donde el arte, las humanidades, la tecnología y las ciencias sociales y exactas, como expresión del conocimiento, hagan parte de la cotidianidad de los colombianos nos exige dejar de actuar (...)
		
	\end{quote}
	
	From the category Environmental sciences and engineering
	
	\begin{quote}
		(...) la justicia ambiental y climática. 4. Articular los objetivos de desarrollo sostenible y la agenda 2030 del posconflicto, focalizándonos en energías renovables, transporte multimodal, bioconstrucción, permacultura, autoconsumo y territorios hídricos resilientes. 5. Redefinir la relación urbanorural, enfatizando restauración y conservación de ecosistemas, equidad, solidaridad, corresponsabilidad e identidad. 6. Revisar el modelo mineroenergético ajustándolo al ordenamiento ambiental territorial, no construir nuevos megaproyectos hidroeléctricos, negar el uso del fracking y hacer transición energética hacia energías renovables alternativas. 7. Promover la educación integral y la democratización de la información y el conocimiento como instrumentos de construcción de la paz, respetando las utopías y las diversidades culturales y naturales. 8. Fortalecer la investigación científica, la innovación y otros modos de construcción del conocimiento como medios para mejorar el conocimiento de la realidad ecológica, climática y cultural. 9. Reconocer a la naturaleza como víctima del conflicto armado y del modelo de desarrollo colombiano, y, por (...)
	\end{quote}
	
	From the category Mathematics and statistics.
	
	\begin{quote}
		(...) la ciencia aspira a los modelos económicos, a reducirlo todo a un cuerpo teórico pequeño de donde se puedan desprender, por deducción directa, sus corolarios básicos. En matemáticas, digamos, es más bella una demostración corta y aritmética que otra más larga que deba recurrir al cálculo para demostrar lo mismo. Todas las geometrías se desarrollan a partir de un puñado de postulados elementales. Los físicos sueñan con la Teoría del Todo están hartos de lidiar con ese enjambre de partículas empalagosas. ¡Eso no es elegante!. En ciencia, como en arte, lo bello es lo breve, lo simple, así haya siempre una secreta complejidad escondida entre los pliegues. Sí, dedicaré lo mejor de mis días a la búsqueda de la fragancia madre y cuando la encuentre la envasaré en un frasquito que será una gema discreta, claro, y la pondré a los pies de mi esposa, lo prometo. El resto será (...)
	\end{quote}
	
	From the category Chemical sciences
	
	\begin{quote}
		(...) a la formación en ciencias de la naturaleza sabemos que el conocimiento científico y tecnológico tiene unos efectos sociales, económicos y políticos que se deben analizar rigurosamente para que nuestros estudiantes fundamenten su formación ciudadana de manera amplia y crítica, de tal manera que cuenten con las herramientas conceptuales y metodológicas para participar activamente en la vida social. Por ejemplo, cuando estudiamos el tema de hidrocarburos, no solo debemos entender su disposición en la naturaleza, sus propiedades y estructuras químicas, sino también analizar su importancia social, los procesos de explotación y los usos sociales que, en su mayoría, están asociados a la generación de energía a partir de la mezcla de hidrocarburos más famosa que conocemos en el mundo y que denominamos petróleo; al estudiar este tema no solo debemos comprenderlo científicamente, sino considerar simultáneamente que esta mezcla de sustancias orgánicas representa el componente central de la matriz energética del (...)
	\end{quote}

\end{document}